\section{Introduction}

Consider the following well-known function which maps over
elements of the list, written in Agda:

\begin{minted}{agda}
map : ∀ {A B : Set} → (A → B) → List A → List B
map f [] = []
map f (x ∷ xs) = f x ∷ map f xs
\end{minted}

Enjoying both composability of this function and
simplicity of its type, we would like to reuse its
definition as much as possible. Up to a certain level, we
would surely be able to do so: after all, its definition
is polymorphic in source and target types, so we can
instantiate it with whatever types we want. But we might
wish for much more; of things that come to our mind are
the following:

\begin{itemize}
    \item Reuse \texttt{map} at higher universe levels,
        e.g., given a \textit{type operator}, map over
        \textit{list of types};
    \item Automagically obtain a proof of the free theorem
        for \texttt{map} -- that is, some term of type
        \begin{multline*}
            \forall A\,A'\,(Q:A\to A'\to\mathtt{Prop}),
            \forall B\,B'\,(R:B\to B'\to\mathtt{Prop}),\\
            \forall f\,f',\left(\forall x\,x',Q\,x\,x'\to
                R(f x)\left(f' x'\right)\right)\to\\
            \forall(\mathtt{xs}:\mathtt{List}\,A)
                (\mathtt{xs}':\mathtt{List}\,A'),
            Q^*\,\mathtt{xs}\,\mathtt{xs}'\to
            R^*(\mathtt{map}\,f\,\mathtt{xs})
                (\mathtt{map}\,f'\,\mathtt{xs}');
        \end{multline*}
    \item Given a group $G$, sets $A$ and $B$ with actions
        $a$, $b$ of $G$ on them, and a function
        $f: A \to B$ which respects $a$ and $b$, prove
        that $\mathtt{map}\;f$ respects actions on
        $\mathtt{List}\,A$ and $\mathtt{List}\,B$ induced
        by $a$ and $b$, respectfully;..
\end{itemize}

The first wish is, of course, granted by adjusting the
definition and making it level-polymorphic. Second one is
possible in a system with \textit{internal parametricity},
as exemplified by Altenkirch and others
\cite{Altenkirch:Parametricity2024}. The last wish,
however, has to be fulfilled manually, albeit with a
simple list induction.

A shift in perspective might help. Instead of tweaking the
definition of \texttt{map}, let's remember all of the
metatheoretic power our type system already has.
Parametricity is a first hint; another is due to Lambek
\cite{Lambek1980} who established that simply-typed lambda
calculus (STLC) is an \textit{internal language} of
\textit{cartesian closed categories} (CCC) \footnote{This
is, of course, part of well-known Curry-Howard-Lambek
correspondence.}. In particular, any term written in STLC
can be interpreted in any cartesian closed category; this
way, you don't have to repeat your construction in any CCC
you meet, you just reinterpret the STLC term you already
have \cite{Elliott:Compcat2017}. This correspondence can
be extended further into dependently typed systems,
establishing correspondence with certain kinds of topoi
and enabling us to do synthetic mathematics: prove general
statements in higher-order intuitionistic logic and
translate them into particular theorems about our topos of
interest \cite{Nawrocki:Lean2026}.

All of this power, unfortunately, lies beyond what usual
type systems can see, since that's their metatheoretical
properties. But what if they were internalized? What if we
had a way to \textit{reinterpret terms of a type system
inside itself}? We experiment with this idea and end up
with two minimal systems, STLC-R (section \ref{STLC-R})
and MLTT-R (section \ref{MLTT-R}), based on simply typed
lambda calculus and Martin-L\"of type theory respectively.
Turns out that their capability of
\textit{self-reinterpretation} is expressive enough to
cover all of the discussed usecases, including restricted
form of effect polymorphism (section
\ref{STLC-R-applications}), synthetic mathematics (section
\ref{MLTT-R-synthetic}), universe polymorphism (section
\ref{MLTT-R-displacement}), parametricity (section
\ref{MLTT-R-parametricity}) and more.

\subsection{Contributions}

\begin{itemize}
    \item We start by giving the syntax and semantics of
        our first type system that reinterprets itself. It
        is based on simply typed lambda calculus and is
        called STLC-R. We prove that STLC-R is sound, and
        that global type inference is decidable (section
        \ref{STLC-R}).
    \item As expected, you can reinterpret STLC-R terms
        in any cartesian closed category you define inside
        it, including categories of differentiable
        manifolds for applications in autodifferentiation,
        Kleisli categories of certain monads and more
        (section \ref{STLC-R-applications}).
    \item Building upon the insights gathered from
        developing STLC-R, we give the syntax and
        semantics of a dependently typed
        self-reinterpreting language based on Martin-L\"of
        type theory called MLTT-R. We prove soundness of
        MLTT-R as well, and motivate bidirectional typing
        discipline for it (section \ref{MLTT-R}).
    \item As we intended, MLTT-R incorporates both a
        restricted form of universe polymorphism and
        external parametricity by virtue of
        reinterpretation only. And, of course, you can use
        MLTT-R as a system for synthetic mathematics with
        internal reuse of synthetic results (section
        \ref{MLTT-R-applications}).
\end{itemize}
